% Generated draft: requires mathematical and source review.
\documentclass[11pt]{article}
\usepackage[margin=25mm]{geometry}
\usepackage{fontspec}
\setmainfont{Latin Modern Roman}
\usepackage{amsmath,amssymb,mathtools,unicode-math}
\setmathfont{Latin Modern Math}
\usepackage{enumitem,needspace,xurl,hyperref}
\setlength{\parindent}{0pt}
\setlength{\parskip}{5pt}
\setlength{\emergencystretch}{3em}
\newcommand{\N}{\mathbb{N}}
\newcommand{\Z}{\mathbb{Z}}
\newcommand{\Q}{\mathbb{Q}}
\newcommand{\R}{\mathbb{R}}
\newcommand{\C}{\mathbb{C}}
\newcommand{\F}{\mathbb{F}}
\newcommand{\A}{\mathbb{A}}
\newcommand{\PP}{\mathbb{P}}
\newcommand{\E}{\mathbb{E}}
\newcommand{\eps}{\varepsilon}
\newcommand{\dd}{\,\mathrm{d}}
\begin{document}
\section*{UnsolvedMath ID: 20002696}
\textbf{Source problem code:} catalogue entry
\subsection{Definitions and mathematical statement}
Let \((M,\tau)\) be a finite von Neumann algebra with faithful normal tracial state and coarse Hilbert \(M\)-bimodule \(\mathcal H=L^2(M)\bar\otimes L^2(M^{\mathrm{op}})\). Interpret the source's closable derivation as defined on an appropriate dense unital \(*\)-algebra core \(\mathcal A\subseteq M\):
\[
 \partial(ab)=a\cdot\partial b+(\partial a)\cdot b.
\]
Write \(\overline\partial\) for its closure and \(\Delta=\overline\partial^{\,*}\overline\partial\geq0\), with the closed-form domain understood. Let \(S_t=e^{-t\Delta}\). Markov/completely-positive hypotheses, if needed for a rigidity interpretation, must be checked rather than inferred from a formal operator sum.

\textbf{Question 36.} Suppose
\[
 a(t)=\sup_{\|x\|_\infty\leq1}\|S_tx-x\|_2\longrightarrow0
 \quad(t\downarrow0).
\]
Must \(\partial\) be inner after extending its target to the algebra of operators affiliated with \(M\bar\otimes M^{\mathrm{op}}\)? Affiliated innerness asks for one affiliated implementer \(\Xi\) with
\[
 \partial x=x\cdot\Xi-\Xi\cdot x\qquad(x\in\mathcal A),
\]
using the canonical extended coarse-bimodule actions and requiring the displayed differences to represent the given values. Hilbert-space innerness additionally requires an implementer in \(\mathcal H\); these are distinct conclusions. The source's literal domain ``\(M\)'' does not specify its unbounded derivation core.

\textbf{Reported spectral equivalence.} Put
\[
 b(R)=\sup_{\|x\|_\infty\leq1}
 \|\mathbf1_{(R,\infty)}(\Delta)x\|_2.
\]
The aid reports that the heat uniformity is equivalent to \(b(R)\to0\) as \(R\to\infty\). For example,
\[
 a(t)\leq tR+b(R),\qquad
 b(R)\leq\frac{a(t)}{1-e^{-tR}}\quad(t,R>0).
\]
This controls spectral tails uniformly on the operator-norm unit ball, not automatically its first energy moment \(\|\overline\partial x\|_2^2\).

\textbf{Reported quantitative sufficient condition.} If the uniform modulus also satisfies the square-Dini condition
\[
 \int_0^1\frac{a(t)^2}{t^2}\,\mathrm dt<\infty,
\]
the aid reports a uniform operator-norm-to-\(\mathcal H\) bound for the closed derivation, with every \(x\in M\) in its form/derivation domain. It then reports an implementer in \(\mathcal H\), hence affiliated innerness as a stronger-than-required conclusion. The square-Dini hypothesis is additional; convergence \(a(t)\to0\) alone is not claimed to imply it.

\textbf{Reported group-cocycle case.} In the standard group-cocycle realization, where \(b(g)\) is a Hilbert cocycle and the heat deformation satisfies
\[
 S_t(u_g)=e^{-t\|b(g)\|^2}u_g,
\]
uniformity on the unit ball forces boundedness of the cocycle. The aid reports that the bounded cocycle is a coboundary, giving innerness in that setting. This argument applies when the stated group/Fourier-multiplier realization and target bimodule hypotheses hold; the source summary does not supply a universal group representation or extend this reduction to arbitrary derivations.

\textbf{Reported atomic boundary example.} The aid describes an explicit weighted atomic family with \(a(t)\to0\), but with no Hilbert-space implementer, while an affiliated implementer exists. Its weights and defining derivation are not included in the supplied summary. This is retained as an existence report requiring verification, not reconstructed numerical data. It shows the distinction between uniform spectral tightness and bounded first energy; it is explicitly not a counterexample to the question's affiliated-innerness conclusion. The bare general implication remains unresolved in the supplied report.
\subsection{Short English statement}
The question asks for affiliated innerness from uniform \(L^2\) heat convergence on the operator unit ball. The aid proves a stronger Hilbert-space conclusion with an extra square-Dini bound and gives a group-cocycle special case. Its atomic example separates Hilbert and affiliated innerness without refuting the original implication.
\subsection{Sources}
\begin{itemize}
\item Local UnsolvedMath registry, permanent ID \(20002696\), catalogue entry, \texttt{aim-probability-notes.json}, \texttt{notes[137]}, \emph{Free analysis}, Question 36, focus group ``Dirichlet Forms, from Classical to Quantum'': \url{https://aimath.org/WWN/freeanalysis/freeanalysis.pdf}.
\item Supplied research summary on spectral tails, square-Dini innerness, group-cocycle derivations and a weighted atomic family. No independent authoritative theorem citations or complete atomic construction are supplied.
\item Attempt 1: partial\_result, status at run ``open'', rigor ``proved'', low novelty confidence; registry status partially\_solved. The machine-generated aid is not peer reviewed and requires independent verification before citation. The supplied verification reports PDF inspection; this draft has not visited it. The trailing extracted 7 is a page artifact.
\end{itemize}
\end{document}
